%File: formatting-instruction.tex
\documentclass[letterpaper]{article}
\usepackage{url,graphicx,xcolor}
\usepackage{booktabs}
\usepackage{multirow}
\usepackage{times}
\usepackage{helvet}
\usepackage{courier}
\usepackage[guidelines]{faikrmod3}
% \usepackage[]{faikrmod3} % without guidelines
\frenchspacing
\setlength{\pdfpagewidth}{8.5in}
\setlength{\pdfpageheight}{11in}



% THE \pdfinfo /Title AND /Author ARE NOT NECESSARY, THEY ARE METADATA FOR THE FINAL PDF FILE
\pdfinfo{
/Title (Insert Your Title Here)
/Author (Put All Your Authors Here, Separated by Commas)}
\setcounter{secnumdepth}{0}  
 \begin{document}
% The file aaai.sty is the style file for AAAI Press 
% proceedings, working notes, and technical reports.
%
\title{Benchmarking continuous optimization techniques for DAG learning}
\author{Daniele Tarek Iaisy\\
Master's Degree in Artificial Intelligence, University of Bologna\\
danieletarek.iaisy@studio.unibo.it
}
\maketitle


\begin{abstract}
\begin{quote}
Learning Bayesian Network (DAG) structures is computationally challenging due to acyclicity
constraints. We compare continuous optimization methods (NOTEARS, DAG-GNN) against classical
algorithms (PC, FGES, Hill Climbing) using synthetic (5-37 nodes) and real-world biological
datasets. Results reveal a distinct trade-off: continuous methods scale better with network
complexity, achieving lower Structural Hamming Distances, whereas classical algorithms scale
better with dataset size but bottleneck on large networks. Notably, the graph-based DAG-GNN
underperformed the simpler NOTEARS baseline in both accuracy and efficiency.
\end{quote}
\end{abstract}


\section{Introduction}

\subsection{Domain}
Learning a Bayesian Network from data can be a difficult task. The space of DAGs over a set of 
nodes is a combinatorial space and traditional approaches have a super-exponential worst case
complexity, mainly owed to the aciclicity constraint to be enforced on the learned DAG.
Continuous numerical optimization techniques can be a viable option to speed up the search.

\subsection{Aim}
The objective of this study is to evaluate the performance of numerical optimization 
algorithms, compared to a baseline of classical algorithm such as PC \cite{spirtes1991pc},
FGES \cite{ramsey2017fges}, and hill climbing \cite{heckerman1995learning}.
We seek to understand how good the reconstructed DAG is and the compute resources needed to 
reconstruct it. We also want to understand the scaling of these resources as the number of 
variables grows or the number of data points grows.

% The main objective of this study is to understand how novel techniques based on numerical 
% optimization work compared against more traditional approaches, both in terms of the fidelity
% of the learned network and in terms of resources necessary to run each training algorithm.\\
% To do so, we investigate the NOTEARS algorithm \cite{zheng2018dags} and the DAG-GNN algorithm
% \cite{yu2019daggnn} compared to the classical algorithms PC \cite{spirtes1991pc},
% FGES \cite{ramsey2017fges}, and hill climbing \cite{heckerman1995learning}. We selected NOTEARS because
% it is the pioneering work that first transformed DAG structure learning from an intractable
% combinatorial search into a purely continuous optimization problem. It works by introducing a
% smooth, exact algebraic characterization of the acyclicity constraint 
% which allows the network structure of a linear structural equation model to be
% efficiently estimated using standard black-box numerical solvers.
% DAG-GNN was subsequently chosen as a natural extension of this continuous framework,
% leveraging the representational power of Graph Neural Networks (GNNs). It works by employing
% a GNN-parameterized Variational Autoencoder (VAE) to capture the underlying data distribution
% alongside a polynomial variant of the continuous acyclicity constraint. We included DAG-GNN
% to investigate another approach in this space on top of NOTEARS and specifically because
% the use of GNN-parametrized VAE seems a natural fit for the task of DAG learning.

\subsection{Method}
The bnlearn repository contains a repository of
prebuilt Bayesian networks that can be used to generate a synthetic dataset, with DAGs ranging from
just 5 variables to 1041. We selected 3 of them in the range of variables 5-37.
To avoid only using synthetic data we also tested on the protein signaling
work introduced in \cite{sachs2005causal} which contains a know ground truth DAG with 11 nodes and a dataset of
11672 real samples. The full experiment covers 4 networks with node count ranging from 5 to 37.\\
We investigate the NOTEARS algorithm \cite{zheng2018dags} and the DAG-GNN algorithm
\cite{yu2019daggnn}. NOTEARS was the first work to transform DAG structure learning
from a combinatorial search into a continuous optimization problem,
by introducing a algebraic characterization of the acyclicity constraint 
which allows the network structure to be estimated using black-box
numerical solvers. DAG-GNN was then chosen as an extension. It works by employing
a GNN-parameterized Variational Autoencoder (VAE) to capture the underlying data distribution,
which seemed a natural fit to learn the DAG structure of the probability distribution.\\
For each experiment, we track Structural Hamming Distance (SHD) between the ground truth DAG
and the reconstructed one together with the execution time of the algorithm.\\

\subsection{Results}
NOTEARS achieved the highest overall performance recovering the best DAGs among all of the tested
algorithms. DAG-GNN performed worse than NOTEARS, while being much more costly to run.
These approaches are not significantly affected by an increase in variable count,
while being much more sensitive to sample count in the 
dataset used, a limitation not observed in classical methods.


\section{Model}
Many of the Bayesian network tested for this study come from the bnlearn repository of networks.
Specifically, we used the CANCER composed of 5 variables, the CHILD composed of 20 variables,
and the ALARM one composed of 37 variables. All of these network include discrete variables
only. For more details on the networks we refer the reader to the bnlearn repository.\\
In addition to these discrete benchmarks, we evaluate algorithms on a real-world continuous
dataset introduced by Sachs et al. \cite{sachs2005causal}, derived from single-cell flow 
cytometry measurements. The ground truth network consists of 11 continuous variables 
representing signaling proteins.


\section{Analysis}

\subsection{Experimental setup}
NOTEARS, DAG-GNN, and FGES are run using the implementations provided by the gCastle
Python toolbox, while PC and Hill Climbing are run using the
implementations provided by the pgmpy library.\\
Each algorithm has its own
hyper-parameters so first we run a hyper-parameter grid search over each dataset for each algorithm
to derive good hyper-parameter values for each dataset-algorithm combination.\\
Then, we execute the algorithm on each dataset and compare the learned DAG against the ground-truth one
computing graph reconstruction metrics, along with execution time of the algorithm. The most important 
metric is SHD, which measures how many edge additions, deletions, or
inversions need to be performed on a graph to transform it into the other (lower means better).\\
For all the networks from the bnlearn repository we generated a
synthetic dataset of 1000 samples, while for the dataset from \cite{sachs2005causal} we used
all 11672 samples as a higher data regime. We also note that in the preliminary phase the
BARLEY graph (48 nodes) from the bnlearn repository
was tested, but ultimately discarded because all of the classical algorithms took extremely long to
train on it.\\
Experiment code is available on Github.\footnote{\url{https://github.com/Tarekun/bayesian-continuous}}

\subsection{Results}
On the very small CANCER NOTEARS and DAG-GNN do not appear to be helpful compared to classical
algorithms as they both recover DAGs with worse SHD than other methods achieved, where fGES is
the clear winner.
As the number of nodes increases, NOTEARS starts showing an edge, achieving lower SHD than any other 
method. Notably, DAG-GNN does somewhat better than
classical algorithms, but does not beat NOTEARS. It is still believable that it could beat
NOTEARS as the network becomes bigger and the relationships between variables more complex, but this
investigation is left as future work. Finally we note that on the biggest network PC was the winner,
beating NOTEARS by 5 points.\\
For execution time, we observe that all algorithms take longer with more variables introduced
in the network, with fGES failing to terminate on ALARM and, together PC, were not appliable to the
BARLEY network (48 nodes) leading to its removal from the experiment.
On the other hand, their scaling w.r.t. the dataset size is much less favourable for numerical algorithms:
Hill climbing, PC, and FGES all take roughly the same amount of time on 
the Sachs et al. dataset that it took them on
the CHILD one; NOTEARS instead took roughly 20 times more time to terminate and DAG-GNN took 16 times.\\
Table 1 contains the performance of each algorithm across all datasets.

\begin{table}[t]
\centering
\footnotesize
\setlength{\tabcolsep}{4pt}
\caption{nSHD is the average "scale normalized" SHD, the SHD divided by the number of edges in the original DAG}
\label{tab:results}
\vspace{4pt}
\begin{tabular}{lrrrrrr}
\toprule
Algorithm & Avg. Time (s) & nSHD & CANCER & Sachs & CHILD & ALARM \\
\midrule
NOTEARS &  44.22 & 0.905 & 0.250 & 0.451 & 0.293 & 0.348 \\
DAG-GNN & 519.09 & 0.993 & 0.112 & 0.195 & 0.137 & 0.161 \\
HC      &   9.62 & 1.282 & 0.215 & 0.237 & 0.333 & 0.274 \\
PC      &   3.17 & 0.928 & 0.374 & 0.477 & 0.442 & 0.455 \\
FGES    &   3.66 & 0.877 & 0.452 & 0.512 & 0.523 & 0.514 \\
\bottomrule
\end{tabular}
\end{table}


\section{Conclusion}
In conclusion, this study reveals a trade-off:
continuous methods yield better structural accuracy for complex networks, while classical
algorithms scale better with larger datasets. Surprisingly, DAG-GNN 
underperformed against simpler baselines. Current limitations
include classical bottlenecks on very large networks and continuous solvers' sensitivity to large
samples. Future work should address these issues by optimizing continuous methods for larger
datasets and evaluating advanced models on more complex graphs. 

% In conclusion, this study highlights a distinct trade-off between continuous optimization and
% classical approaches in Bayesian network structure learning, revealing that continuous methods
% excel in structural accuracy as network complexity grows, whereas classical algorithms demonstrate
% superior scalability with respect to dataset size. An interesting and somewhat unexpected finding
% was the underperformance of DAG-GNN; despite its architecturally natural fit for graph-structured 
% data, it failed to surpass the simpler continuous baseline while incurring significantly higher
% computational costs. A primary limitation of our current experimental setup is the computational
% bottleneck experienced by classical algorithms on very large networks, which restricted evaluation
% to moderately sized graphs, alongside the observed sensitivity of continuous solvers to large
% sample sizes. Future work should therefore focus on optimizing continuous methods for larger
% datasets and evaluating these advanced graph-based models on substantially more complex networks
% to determine if their representational advantages ultimately manifest at a much larger scale.


\section{Links to external resources}
\begin{itemize}
    \item pgmpy library: \url{https://pgmpy.org/index.html}
    \item gCastle library: \url{https://github.com/huawei-noah/trustworthyAI/blob/master/gcastle/README.md}
    \item bnlearn repository: \url{https://www.bnlearn.com/bnrepository/}
\end{itemize}

\bibliographystyle{aaai}
\bibliography{faikrmod3.bib}


\end{document}